\chapter{Родитель \texorpdfstring{$\Delta$}{Delta} moment-map odd
auxiliary trilinear}
\label{ch:v10-h15-r2-delta-moment-map-odd-trilinear-parent}

Предыдущий аудит выделил член
\begin{equation}
 \langle A_\Sigma,Q\Sigma\rangle,
 \qquad Q=6Y,
\end{equation}
как ближайший единый источник смешанного блока. Проверим, возникает ли он
из канонического отображения момента на общей паре изоморфных модулей
$A_\Sigma\oplus\Sigma$.

Нейтральный $\Delta$-фон действительно фиксирует направление оператора:
\begin{equation}
 Q=6\mu_R-4\operatorname{ad}(\mu_4)
 =\operatorname{diag}(3,-3,7,1,-1,-7,3,-3).
\end{equation}
Но отображение момента ортогональной прямой суммы аддитивно. При
канонической метрике кратности $K_0=I_2$ соответствующий оператор равен
$K_0\otimes Q$ и его $A_\Sigma$--$\Sigma$ блок имеет ранг нуль.

Общая Real-симметричная инвариантная метрика двух изоморфных копий после
единичной нормировки диагоналей имеет вид
\begin{equation}
 K(\kappa)=
 \begin{pmatrix}1&\kappa\\\kappa&1\end{pmatrix}.
\end{equation}
Нормировка фиксирует два из трёх симметричных коэффициентов и оставляет
$\kappa$ свободным. Mixed moment-map block становится $\kappa Q$.
Например, $\kappa=1/2$ даёт невырожденную положительную метрику, но лишь
$Q/2$. Точное значение $Q$ требует $\kappa=1$.

Hadamard congruence даёт
\begin{equation}
 H^TK(0)H=\operatorname{diag}(2,2),\qquad
 H^TK(1/2)H=\operatorname{diag}(3,1),\qquad
 H^TK(1)H=\operatorname{diag}(4,0).
\end{equation}
Следовательно, target-значение имеет rank $1$ на пространстве кратности и
rank/nullity $8/8$ на полном 16-мерном trace carrier: комбинация
$A_\Sigma-\Sigma$ теряет норму, и независимое auxiliary-поле исчезает.
Антисимметричное смешивание не помогает, поскольку оно задаёт
кососимметричный оператор и не вносит Real quadratic scalar.

\begin{theorem}[Аддитивность moment map и свободная метрика кратности]
$\Delta$ moment map единственно фиксирует Cartan-направление $Q$, но
каноническая ортогональная сумма даёт нулевой cross-блок. Инвариантное
смешивание порождает лишь семейство $\kappa Q$; нормировка не определяет
$\kappa$, а точное значение $\kappa=1$ вырождает trace metric и уничтожает
независимость auxiliary-копии.
\end{theorem}

\begin{nogocriterion}[Вырожденная multiplicity metric не является parent-origin]
Нельзя выводить $\langle A_\Sigma,Q\Sigma\rangle$ из отображения момента,
выбирая off-diagonal pairing после просмотра target. При $\kappa<1$
коэффициент свободен и неверен, при $\kappa=1$ одна полевая комбинация
имеет нулевую норму, а унаследованное значение равно $\kappa=0$.
\end{nogocriterion}

\begin{protocol}\begin{enumerate}
\item exact moment-map generator rank: \textbf{$8$};
\item symmetric invariant metric basis rank: \textbf{$3$};
\item unit-normalization rank/nullity: \textbf{$2/1$};
\item inherited/required mixing: \textbf{$\kappa=0/1$};
\item inherited/required cross rank: \textbf{$0/8$};
\item nondegenerate test $\kappa=1/2$: \textbf{$Q/2$};
\item target multiplicity metric rank/nullity: \textbf{$1/1$};
\item target full trace metric rank/nullity: \textbf{$8/8$};
\item conditional/physical origin: \textbf{$3/5$ и $3/5$};
\item следующий гейт: \textbf{родитель mixed superconnection curvature}.
\end{enumerate}\end{protocol}

Машинный сертификат сохранён в
\path{s2t_v10_particle_wrinkle_dislocation_callias_equal_charge_twist_m4_h15_r2_delta_moment_map_odd_auxiliary_trilinear_parent_origin_gate_results.json}.